\documentclass[11pt,a4paper]{article}
\usepackage[a4paper,margin=24mm]{geometry}
\usepackage[T1]{fontenc}
\usepackage[utf8]{inputenc}
\usepackage{lmodern}
\usepackage{amsmath,amssymb,amsthm}
\usepackage{graphicx}
\usepackage{booktabs,array,tabularx,multirow}
\usepackage[font=small,labelfont=bf]{caption}
\usepackage{enumitem}
\usepackage{microtype}
\usepackage{xcolor}
\definecolor{linkblue}{rgb}{0,0.15,0.5}
\usepackage[colorlinks=true,linkcolor=linkblue,citecolor=linkblue,urlcolor=linkblue]{hyperref}
\usepackage{url}
\graphicspath{{./}{figs/}{paper11c_worked_systems_audit_v5_src/figs/}}
\theoremstyle{plain}
\newtheorem{theorem}{Theorem}[section]
\newtheorem{proposition}[theorem]{Proposition}
\newtheorem{lemma}[theorem]{Lemma}
\newtheorem{corollary}[theorem]{Corollary}
\theoremstyle{definition}
\newtheorem{definition}[theorem]{Definition}
\newtheorem{example}[theorem]{Example}
\newtheorem{convention}[theorem]{Convention}
\theoremstyle{remark}
\newtheorem{remark}[theorem]{Remark}
\newcommand{\R}{\mathbb{R}}
\newcommand{\Qr}{\mathbb{Q}}
\newcommand{\N}{\mathbb{N}}
\newcommand{\calA}{\mathcal{A}}
\newcommand{\calB}{\mathcal{B}}
\newcommand{\calF}{\mathcal{F}}
\newcommand{\calK}{\mathcal{K}}
\newcommand{\calN}{\mathcal{N}}
\newcommand{\calP}{\mathcal{P}}
\newcommand{\calS}{\mathcal{S}}
\newcommand{\calW}{\mathcal{W}}
\newcommand{\Vstar}{V^{\ast}}
\newcommand{\Rstar}{R^{\ast}}
\newcommand{\Gstar}{G^{\ast}}
\newcommand{\tfr}[2]{\tfrac{#1}{#2}}
\newcommand{\st}[1]{\mathsf{#1}}
\newcommand{\led}{\ell}
\newcommand{\priv}{\pi}
\newcommand{\opt}{\ell^{+}}
\newcommand{\free}{\mathrm{f}}
\newcommand{\nr}{\mathrm{nr}}
\newcommand{\ofull}{o_{\mathrm{id}}}
\newcommand{\oagg}{o_{\Sigma}}
\newcommand{\oblind}{o_{\mathrm{b}}}
\newcommand{\Kinst}{K_{\mathrm{inst}}}
\newcommand{\Kagg}{K_{\mathrm{agg}}}
\newcommand{\Kcod}{K_{\mathrm{cod}}}
\newcommand{\Kfull}{K_{\mathrm{full}}}
\newcommand{\Kdec}{K_{\mathrm{dec}}}
\newcommand{\NCHECKS}{64}
\setlength{\parskip}{2pt}
\emergencystretch=2em
\allowdisplaybreaks

\begin{document}

\title{\textbf{Worked systems for viability under incomplete observation:\\ exact belief kernels across observation structures, protocols and decentralisation,\\ a monitoring-design census, and certified margins}}
\author{Amin Abaee\\[0.3em]
{\small Independent Researcher}\\[0.2em]
{\small ORCID: \href{https://orcid.org/0000-0002-0019-1842}{0000-0002-0019-1842}\quad
\href{mailto:amin_abaee@ut.ac.ir}{amin\_abaee@ut.ac.ir}}}
\date{October 15, 2026}
\maketitle

\begin{abstract}
\noindent
A viability problem under incomplete observation is specified by much more than a dynamics and a safe set: the observation map, the protocol that constrains which instruments may follow which, the convention by which a shared record is kept, and the allocation of authority between agencies all enter the answer. This article develops an exact \emph{audit} of these choices on a small worked family and extracts the general theorems that the computations expose. The worked system has two regenerating stocks on the grid $\{0,\dots,3\}^2$, two instruments that each regenerate one stock while the other declines, and the floor $z_1,z_2\ge 1$; beliefs are sets of states, and viability is decided by an exact fibrewise belief game whose greatest fixed point is computed over all $511$ beliefs. A single monotonicity theorem orders the kernels along three axes---refinement of the observation, relaxation of the protocol, and horizon---and the worked system shows every axis strict and the two middle classes incomparable ($111$, $143$, $119$ and $255$ viable beliefs for the institutional, aggregate, codex and full classes). With two instruments the no-repeat protocol is strict alternation, and the two conventions for the register admit closed forms: the shared-ledger kernel depends on the observation only at the first step, while the private-register kernel is a phase-coherence condition; the conventions are incomparable, and the aggregate kernel is realised by the $32$ alternation policies although no memoryless policy sustains a single state. Decentralised unanimity factors statewise, so a belief is sustained if and only if its states share a sustaining law pair: $12$ of the $28$ coordinated-viable pairs survive, and the allocation of one vote moves the kernel between $12$ and $4$ pairs. Monitoring design is reduced to graph structure: one-step adequacy of a partition is adequacy against an incompatibility graph, and the whole set of individually viable states is dynamically viable under a partition if and only if the partition is one-step adequate on it; the census over all $4{,}140$ partitions of the eight viable states finds $1{,}694$ adequate designs, $16$ coarsest two-cell designs, $74$ distinct kernels, and no adequate design at all under the shared ledger. Review timing on regime graphs is a max-plus hitting time with $33$ and $43$ viable cells on a $93$-cell grid, and the continuous companion fibre is certified exactly (crossover $Y^{\ast}=27/5$, margin law $(Y-27/5)/10$, saddle value $-1/10$, Helly-tight triple). Every number is regenerated by an exact script with two independent implementations of the kernels; the monotonicity theorem is formalised in Lean~4.
\end{abstract}

\medskip
\noindent\textbf{Keywords:} viability kernel; partial observation; belief games; information structures; protocols and memory; decentralised control; monitoring design; set partitions; max-plus algebra; LP duality; exact computation; Lean~4

\medskip
\noindent\textbf{MSC 2020:} 93B03, 93A14, 91A43, 05A18, 90C05, 68Q60, 15A80
